\subsection{THE MULTIPLICATIVE GROUP R*}

We know from Chapter III, § 6, no. 7, that the topology induced on \(\mathbf{R}^*\) by the topology of the real line is compatible with the multiplicative group structure of \(\mathbf{R}^*\) ; since \(\mathbf{R}^*\) is an open subset of the locally compact space \(\mathbf{R}\) , it follows that \(\mathbf{R}^*\) is a locally compact topological group (Chapter I, § 9, no. 7, Proposition 13) and is therefore complete (Chapter III, § 3, no. 3, Corollary I to Proposition 4; this follows also from Chapter III, § 6, no. 8, Proposition 8); of course, this latter property relates to the multiplicative uniformity on \(\mathbf{R}^*\) and not to the uniformity induced on \(\mathbf{R}^*\) by the additive uniformity of \(\mathbf{R}\) .

The function \(xy\) maps the set \(\mathbf{Q}_{+} \times \mathbf{Q}_{+}\) into \(\mathbf{Q}_{+}\) , and therefore it maps \(\mathbf{R}_{+} \times \mathbf{R}_{+}\) into \(\mathbf{R}_{+}\) (Chapter I, § 2, no. 1, Theorem 1); in other words, the product of two real numbers \(\geqslant 0\) is \(\geqslant 0\) . The formulae \((-x)y = -xy\) and \((-x)(-y) = xy\) then show that the product of a number \(\geqslant 0\) and a number \(\leqslant 0\) is \(\leqslant 0\) , and that the product of two numbers \(\leqslant 0\) is \(\geqslant 0\) ; from this it follows that

\[
| x y | = | x |. | y |\tag{1}
\]

(which might also have been obtained by extension of the corresponding relation on \(\mathbf{Q} \times \mathbf{Q}\) ).

If \(x > 0\) and \(y > 0\) we have \(xy \neq 0\) , and therefore \(xy > 0\) ; likewise, if \(x < 0\) and \(y > 0\) , then \(xy < 0\) ; and if \(x < 0\) and \(y < 0\) , then \(xy > 0\) . In particular, if \(x \neq 0\) we have \(x^2 > 0\) , so that a sum of squares of real numbers cannot be zero unless each of the numbers is zero.

If \(x > 0\) and \(y \leqslant z\) (resp. \(y < z\) ) we have \(xy \leqslant xz\) (resp. \(xy < xz\) ) in other words, a homothety of ratio \(> 0\) preserves order on \(\mathbf{R}\) . Since \((-x)y = -xy\) , a homothety of ratio \(< 0\) changes the ordering on \(\mathbf{R}\) into the opposite ordering.

If \(x > 0\) we have \(\mathbf{1} / x > 0\) , since \(x.(\mathbf{1} / x) = \mathbf{1} > 0\) . If \(0 < x < y\) we have \(xy > 0\) , hence \(x.(\mathbf{1} / xy) < y.(\mathbf{1} / xy)\) , that is \(\mathbf{1} / y < \mathbf{1} / x\) . Hence the mapping \(x \to \mathbf{1} / x\) of the set \(\mathbb{R}_+^*\) of real numbers \(> 0\) onto itself is strictly decreasing.

We see in the same way that the function \(\mathbf{1} / x\) is strictly decreasing in \(] \leftarrow, o[, \text{ and therefore the function } \frac{\mathbf{1}}{x - a}\) is strictly decreasing in each of the intervals \(] \leftarrow, a[ \text{ and } ]a, \rightarrow [.\)

It follows from what precedes that \(R_{+}^{*}\) is a subgroup of the multiplicative group \(R^{*}\) ; moreover, the order relation \(x \leqslant y\) is compatible with the multiplicative, group structure of \(R_{+}^{*}\) ; in other words, \(R_{+}^{*}\) is a linearly ordered group.

The fact that the product of two real numbers \(\geqslant0\) is \(\geqslant0\) can be expressed by saying that R is an ordered field: all the above properties are common to all ordered fields.

PROPOSITION 2. The multiplicative group \(\mathbf{R}^{*}\) of real numbers \(\neq 0\) is a topological group isomorphic to the product of its subgroups \(\mathbf{R}_{+}^{*}\) and \(\mathbf{U}_{0}\) , where

\[
\mathrm{U} _ {0} = \{- 1, + 1 \}.
\]

For each \(x \neq 0\) let \(\operatorname{sgn} x\) denote \(\frac{x}{|x|}\) (sign of \(x\) ). The function \(\operatorname{sgn}\) is a homomorphism of \(\mathbf{R}^*\) onto \(U_0\) . We have \(x = |x| \operatorname{sgn} x\) , and this decomposition of \(x\) as the product of an element of \(\mathbf{R}_+^*\) and an element of \(U_0\) is unique; hence the group structure of \(\mathbf{R}^*\) is the product of the group structures of \(\mathbf{R}_+^*\) and \(U_0\) . On the other hand, the mapping \(x \to |x|\) is continuous, and so is \(x \to \operatorname{sgn} x = \frac{x}{|x|}\) , since \(x \neq 0\) . Hence the result.

We extend the function sgn to the whole of R by putting sgn o = o.

We shall see in Chapter V (§ 4, no. 1, Theorem 1) that the topological group \(\mathbf{R}_{+}^{*}\) is isomorphic to the additive group \(\mathbf{R}\) ; this will complete the determination of the structure of the topological group \(\mathbf{R}^{*}\) .

